% TODO make hw stylesheet
\documentclass[10pt]{article}
\usepackage[letterpaper,top=1in,left=1in,right=1in]{geometry}
\usepackage[dvipsnames,svgnames]{xcolor}
\usepackage[shortlabels]{enumitem}
\usepackage[framemethod=TikZ]{mdframed}
\usepackage{amsmath,amssymb,amsthm}
\usepackage[colorlinks]{hyperref}
\usepackage{multicol}
\usepackage{tabularx}
\usepackage{bbm}

\hypersetup{
    colorlinks=true,
    linkcolor=blue,
    filecolor=magenta,
    urlcolor=magenta,
    pdftitle={MAT 344 HW}
}

\usepackage{mathtools}
\usepackage{thmtools}
\usepackage[textsize=footnotesize]{todonotes}
\usepackage{titling}
\usepackage{cancel}

\reversemarginpar
\mdfdefinestyle{mdbluebox}{roundcorner=10pt,innerbottommargin=9pt,
linecolor=blue,backgroundcolor=TealBlue!5,}
\declaretheoremstyle[headfont=\sffamily\bfseries\color{MidnightBlue},
mdframed={style=mdbluebox},]{thmbluebox}

\mdfdefinestyle{mdbloobox}{frametitlefont=\bfseries,innerbottommargin=8pt,
nobreak=true,backgroundcolor=TealBlue!5,linecolor=blue,}
\declaretheoremstyle[headfont=\bfseries\color{MidnightBlue},
mdframed={style=mdbloobox},headpunct={\\[3pt]},postheadspace=0pt,]{thmbloobox}

\mdfdefinestyle{mdredbox}{frametitlefont=\bfseries,innerbottommargin=8pt,
nobreak=true,backgroundcolor=Salmon!5,linecolor=RawSienna,}
\declaretheoremstyle[headfont=\bfseries\color{RawSienna},
mdframed={style=mdredbox},headpunct={\\[3pt]},postheadspace=0pt,]{thmredbox}

\mdfdefinestyle{mdgreenbox}{linecolor=ForestGreen,backgroundcolor=ForestGreen!5,
linewidth=2pt,rightline=false,leftline=true,topline=false,bottomline=false,}
\declaretheoremstyle[headfont=\bfseries\sffamily\color{ForestGreen!70!black},
mdframed={style=mdgreenbox},headpunct={ --- },]{thmgreenbox}

\mdfdefinestyle{mdorangebox}{linecolor=RedOrange,backgroundcolor=RedOrange!5,
linewidth=2pt,rightline=false,leftline=true,topline=false,bottomline=false,}
\declaretheoremstyle[headfont=\bfseries\sffamily\color{RedOrange!70!black},
mdframed={style=mdorangebox},headpunct={ --- },]{thmorangebox}

\mdfdefinestyle{mdblackbox}{linecolor=black,backgroundcolor=RedViolet!5!gray!5,
linewidth=3pt,rightline=false,leftline=true,topline=false,bottomline=false,}
\declaretheoremstyle[mdframed={style=mdblackbox}]{thmblackbox}

\declaretheoremstyle[spaceabove=6pt,spacebelow=6pt]{boldhead}
\declaretheoremstyle[spaceabove=6pt,spacebelow=6pt,headfont=\normalfont\itshape]{italichead}

\declaretheorem[style=boldhead,name=Problem]{problem}
\declaretheorem[style=boldhead,name=Problem,numbered=no]{problem*}
\declaretheorem[style=boldhead,name=Setup]{setup} % for config geo
\declaretheorem[style=boldhead,name=Example,sibling=problem]{example}
\declaretheorem[style=boldhead,name=Theorem,sibling=problem]{theorem}
\declaretheorem[style=boldhead,name=Corollary,sibling=theorem]{corollary}
\declaretheorem[style=boldhead,name=Proposition,sibling=theorem]{prop}
\declaretheorem[style=boldhead,name=Lemma,numberwithin=problem]{lemma}
\declaretheorem[style=boldhead,name=Theorem,numbered=no]{theorem*}
\declaretheorem[style=boldhead,name=Lemma,numbered=no]{lemma*}
\declaretheorem[style=boldhead,name=Claim,numberwithin=problem]{claim}
\declaretheorem[style=boldhead,name=Claim,numbered=no]{claim*}
\declaretheorem[style=boldhead,name=Remark,numberwithin=problem]{remark}
\declaretheorem[style=boldhead,name=Remark,numbered=no]{remark*}
\declaretheorem[style=boldhead,name=Definition,sibling=theorem]{definition}
\declaretheorem[style=boldhead,name=Definition,numbered=no]{definition*}

\providecommand{\ol}{\overline}
\providecommand{\ceil}[1]{\left\lceil #1 \right\rceil}
\providecommand{\floor}[1]{\left\lfloor #1 \right\rfloor}
\newcommand{\dd}{\mathrm{d}}
\providecommand{\ul}{\underline}
\providecommand{\eps}{\varepsilon}
\providecommand{\half}{\frac{1}{2}}
\providecommand{\CC}{\mathbb C}
\providecommand{\FF}{\mathbb F}
\providecommand{\II}{\mathbb I}
\providecommand{\NN}{\mathbb N}
\providecommand{\QQ}{\mathbb Q}
\providecommand{\RR}{\mathbb R}
\providecommand{\ZZ}{\mathbb Z}
\providecommand{\ind}{\mathbbm 1}
\providecommand{\dg}{^\circ}
\providecommand{\ii}{\item}
\providecommand{\setdiff}{\smallsetminus}
\providecommand{\alert}[1]{{\sffamily\textbf{\textcolor{blue}{#1}}}}
\providecommand{\inv}{^{-1}}
\providecommand{\mb}[1]{\mathbf{#1}}
\newcommand{\what}{\widehat}

\newenvironment{soln}{\begin{proof}[Solution]}{\end{proof}}

\providecommand{\pow}{\operatorname{Pow}}
\providecommand{\vocab}[1]{{\textbf{\textcolor{ForestGreen}{#1}}}}
\providecommand{\lcm}{\operatorname{lcm}}
\providecommand{\abs}[1]{\left\lvert #1\right\rvert}
\providecommand{\smabs}[1]{\lvert #1\rvert}
\providecommand{\innerprod}[1]{\langle\!\langle #1\rangle\!\rangle}
\providecommand{\norm}[1]{\left\| #1\right\|}
\providecommand{\smnorm}[1]{\| #1\|}
\providecommand{\gr}{\operatorname{gr}}
\providecommand{\im}{\operatorname{im}}
\providecommand{\paren}[1]{\left(#1\right)}

\usepackage{mathrsfs}

% place title at top right
\setlength{\droptitle}{-8ex}
\pretitle{\begin{flushright}\Large\bfseries}
\posttitle{\par\end{flushright}}
\preauthor{\begin{flushright}}
\postauthor{\end{flushright}}
\predate{\begin{flushright}}
\postdate{\end{flushright}}

\title{MAT 344 HW07}
\author{\large Aditya Pahuja}
\date{\large Due 04/07}
\begin{document}
\maketitle

\begin{problem}
    [p.171\#5]
    Show that the mean value formula remains valid for
    $u = \log\abs{1+z}$, $z_0 = 0$, $r = 1$, and use this fact
    to compute
    \begin{equation*}
        \int_0^\pi\log\sin\theta\dd\theta.
    \end{equation*}
\end{problem}

\begin{soln}
    Observe that $\log\abs{1 + z}$ is harmonic in
    the open disk $B_1(0)$ because it is the real part
    of the holomorphic function $\log(1 + z)$
    defined in the simply connected domain $B_1(0)$,
    so certainly
    \begin{equation*}
	0 = u(0) = \frac 1{2\pi} \int_0^{2\pi}
	\log\abs{1 + re^{i\theta}}\dd\theta
    \end{equation*}
    for $0 \le r < 1$. We need to show that
    \begin{equation*}
	\lim_{r\to 1^-} \frac 1{2\pi} \int_0^{2\pi}
	\log\abs{1 + re^{i\theta}}\dd\theta
	= \frac 1{2\pi} \int_0^{2\pi}
	\log\abs{1 + e^{i\theta}}\dd\theta.
    \end{equation*}
    Consider the function $g\colon[0,2\pi]\to\RR$ given by
    \begin{equation*}
	g(\theta) = \max\paren{1, \frac 2{\sqrt{\abs{\pi-\theta}}}}.
    \end{equation*}
    Clearly $g$ is integrable on $[0,2\pi]$
    because $\frac 1{\sqrt x}$ is integrable on any bounded subset of $\RR$.
    Moreover, we can bound
    \begin{equation*}
	\frac{\abs{\pi - \theta}}2 <
	\abs{\sin\theta} \le \abs{1 + e^{i\theta}}
    \end{equation*}
    for $\theta$ sufficiently close to $\pi$
    (since $\abs{\sin\theta}$ is the distance from $1 + e^{i\theta}$
    to the real axis), so
    \begin{equation*}
	\abs{\log(1 + e^{i\theta})}
	\le \frac{1}{\sqrt{\abs{1 + e^{i\theta}}}}
	\le \frac 2{\sqrt{\abs{\pi - \theta}}}.
    \end{equation*}
    We can note that for a given $\theta$,
    $\abs{1 + re^{i\theta}}$ is maximized over $r\in[0,1]$
    at either $r = 0$ or $r = 1$
    because the squared distance is quadratic in $r$,
    hence is a convex function of $r$, so
    $g(\theta)$ is an upper bound for $\abs{\log\abs{1 + re^{i\theta}}}$.
    for all $r\in[0,1]$ and $\theta\in[0,2\pi]$,
    hence the dominated convergence theorem implies
    the desired convergence.

    Using the identity $1 + \cos\theta = 2\cos^2(\frac\theta2)$
    and the integral computed above,
    \begin{align*}
	0 =\int_0^{2\pi} \log\smabs{1 + e^{i\theta}}\dd\theta
	&= \int_0^{2\pi}\log(\sqrt{2 + 2\cos\theta})\dd\theta\\
	&= \int_0^{2\pi}\log
	\paren{2\abs{\cos\paren{\frac\theta2}}}\dd\theta
	= 2\int_0^{\pi}\log\paren{2\cos\frac\theta2}
	\dd\theta\\
	&= \pi\log 2 +
	\int_0^\pi\log\paren{\cos\frac\theta2}\dd\theta
	= \pi\log 2 + \int_0^{\pi}\log\paren{\sin\frac\theta2}\dd\theta\\
	&= \pi\log2 + 2\int_0^{\pi/2}\log(\sin u)\dd u
	= \pi\log2 + \int_0^{\pi}\log(\sin u)\dd u,
    \end{align*}
    so the integral in the last line,
    which is what we want to find, is equal to $\boxed{-\pi\log 2}$.
\end{soln}

\begin{problem}
    [p.171\#6]
    If $f(z)$ is analytic in the whole plane and if
    $z\inv\operatorname{Re}f(z)\to 0$ when $z\to\infty$,
    show that $f$ is a constant.
\end{problem}

\begin{soln}
    Let $u$ be the real part of $g(z) = f(z) - f(0)$,
    so that equation (66) says
    \begin{equation*}
	g(z) = \frac 1{2\pi i}\int_{\abs\zeta = R}
	\frac{\zeta + z}{\zeta - z}\cdot u(\zeta)\frac{\dd\zeta}{\zeta}
	+ iC
    \end{equation*}
    for some constant $C$, if $\abs z < R$.
    We have that $\lim_{z\to \infty}\frac{u(z)}{z} = 0$
    and $\lim_{z\to 0} \frac{g(z)}{z} = f'(0)$,
    implying that $\frac{u(z)}{z}$ is bounded everywhere,
    so suppose that $\abs{u(z)} \le M\abs z$ for some constant $M > 0$.
    Taking $R = 2\abs z$, we find
    \begin{equation*}
	\abs{f(z)} \le \frac 1{2\pi}
	C + \int_{\abs\zeta = R}
	\abs{\frac{Re^{i\theta} + z}{Re^{i\theta} - z}}
	u(Re^{i\theta})\dd\theta
	\le C + 3MR = C + 6M\abs z,
    \end{equation*}
    noting that $\abs{Re^{i\theta} + z} \le 3\abs z$
    and $\abs{Re^{i\theta} - z} \ge \abs z$.
    We showed in a previous homework that this implies
    $f$ is a linear polynomial $f(z) = az + b$. Then, if $a\ne 0$,
    \begin{equation*}
	a
	= \lim_{\substack{r\to\infty\\r\in\RR}} \frac{r\abs a^2 + b}{r\ol a}
	= \lim_{\substack{r\to\infty\\r\in\RR}} \frac{u(r\ol a) +
	\operatorname{Im}b}{r\ol a}
	= \lim_{z\to\infty} \frac{u(z)}{z},
    \end{equation*}
    so the right-hand limit is nonzero,
    which contradicts the given information;
    thus $a$ is zero and $f(z)\equiv b$.
\end{soln}

\begin{problem}
    [p.174\#3]
    If $f(z)$ is analytic in $\abs z\le 1$ and satisfies
    $\abs f = 1$ on $\abs z = 1$, show that $f(z)$ is rational.
\end{problem}

\newcommand{\HH}{\mathbb H}
\begin{soln}
    Let $L$ be a fractional linear transformation
    biholomorphically mapping the closed disk $\ol\Delta$
    to the closed upper half plane $\ol{\mathbb H}$
    (where $\mathbb H = \{z\in\CC : \operatorname{Im}z > 0\}$
    and the closure is taken in the Riemann sphere,
    so that $\partial\HH = \RR\cup\infty$),
    such as $L(z) = i\frac{1+z}{1-z}$.
    Then, consider $g = LfL\inv\colon\ol{\mathbb H}\to\ol{\mathbb H}$.
    Assuming $f$ is nonconstant,
    the maximum principle tells us that
    $f(\Delta)\subseteq\Delta$
    and we are given that $f(\partial\Delta)\subseteq\partial\Delta$.
    This translates to $g(\HH) \subseteq\HH$
    and $g(\partial\HH)\subseteq\partial\HH$,
    meaning in particular that the poles and zeroes of $g$
    are all real; thus
    \begin{equation*}
	g(z) = \frac{(z-z_1)\cdots(z-z_n)}{(z-w_1)\cdots(z-w_m)}h(z)
    \end{equation*}
    for $z_i,w_i\in\RR$ and $h\in\mathcal O(\ol\HH)$
    a nowhere zero and nowhere singular holomorphic function
    with $h(\partial \HH)\subseteq\RR - \{0\}$.

    I claim that $h$ must be a constant function.
    Let $k =  L\inv hL\colon\ol\Delta\to\ol\Delta$;
    since $h(\partial\HH)\subseteq\partial\HH - \{0\}$,
    $k(\partial\Delta)\subseteq\partial\Delta - \{L\inv(0) = -1\}$.
    It is easy to see (e.g. by stereographic projection)
    that $\partial\Delta - \{-1\}$ is homeomorphic to
    the contractible space $\RR$, so
    $k|_{\partial\Delta}\colon\partial\Delta\to\partial\Delta$
    is nullhomotopic, meaning the winding number of
    $k$ around any interior point of $\partial\Delta$ is zero.
    By the argument principle, we therefore find that
    $k(\Delta)\subseteq\partial\Delta$,
    so the fact that nonconstant holomorphic functions
    are open maps tells us that $k$, hence $h$, is constant;
    in fact $h(z)\equiv c$ for some nonzero real $c$.

    We conclude that
    \begin{equation*}
        LfL\inv(z) = g(z) = 
	c\cdot\frac{(z-z_1)\cdots(z-z_n)}{(z-w_1)\cdots(z-w_m)}
    \end{equation*}
    is rational, so certainly the composition of rational functions
    $g = L\inv fL$ is also rational.
\end{soln}

\pagebreak
\begin{problem}
    [p.174\#4]
    Use (66) to derive a formula for $f'(z)$ in terms of $u(z)$.
\end{problem}

\begin{soln}
    Equation (66) says that for $z,w\in\CC$
    with $\abs z,\abs w < R$,
    \begin{align*}
	\frac{f(z)-f(w)}{z-w}
	&= \frac 1{z-w}\cdot\frac 1{2\pi i}\int_{\abs\zeta = R}
	\paren{\frac{\zeta+z}{\zeta-z}
	- \frac{\zeta+w}{\zeta-w}}\cdot u(\zeta)\frac{\dd\zeta}{\zeta}\\
	&= \frac 1{2\pi i}\int_{\abs\zeta = R}
	\frac 1{z-w}\cdot\frac{2(z-w)\zeta}{(\zeta-z)(\zeta-w)}
	\cdot u(\zeta)\frac{\dd\zeta}{\zeta}
	= \frac 1{2\pi i}\int_{\abs\zeta = R}
	\frac{2u(\zeta)}{(\zeta-z)(\zeta-w)}\dd\zeta.
    \end{align*}
    Take $\delta > 0$ such that $\abs w + \delta < R$
    and $z$ such that $\abs{w-z} < \delta$.
    The magnitude of the expression
    \begin{equation*}
	\frac 1{2\pi i}\int_{\abs\zeta = R}
	\frac{2u(\zeta)}{(\zeta-z)(\zeta-w)}
	-\frac{2u(\zeta)}{(\zeta-w)^2}\dd\zeta
	= \frac 1{2\pi i}\int_{\abs\zeta = R}
	\frac{2u(\zeta)}{\zeta-w}\cdot
	\frac{z-w}{(\zeta-z)(\zeta-w)}\dd\zeta
    \end{equation*}
    can be bounded above by
    \begin{equation*}
	\frac 1{2\pi}\int_0^{2\pi}\abs{\frac{2u(\zeta)}{\zeta-w}}
	\cdot \abs{\frac{z-w}{(\zeta-z)(\zeta-w)}}\cdot R\dd\theta
    \end{equation*}
    where of course $\zeta = Re^{i\theta}$.
    Since $\abs{\zeta-a} \ge R - \abs a$
    for any $a\in\CC$ (noting that $\abs\zeta = R$)
    and $\abs z$ and $\abs w$ are both bounded above by
    $\abs w + \delta$, we have
    \begin{equation*}
        \abs{\frac{z-w}{(\zeta-z)(\zeta-w)}}
	< \frac{\delta}{(R - \abs z)(R - \abs w)}
	< \frac{\delta}{(R - \abs w - \delta)^2},
    \end{equation*}
    so
    \begin{equation*}
	\frac 1{2\pi}\int_0^{2\pi}\abs{\frac{2u(\zeta)}{(\zeta-w)}}
	\cdot \abs{\frac{z-w}{(\zeta-z)(\zeta-w)}}\cdot R\dd\theta
	< \frac{\delta}{(R - \abs w - \delta)^2}
	\cdot \frac R{2\pi}\int_0^{2\pi}
	\abs{\frac{u(\zeta)}{\zeta-w}}\dd\theta.
    \end{equation*}
    The final expression goes to zero with $\delta$, meaning
    \begin{equation*}
	\lim_{z\to w}\frac 1{2\pi i}\int_{\abs\zeta = R}
	\frac{2u(\zeta)}{(\zeta-z)(\zeta-w)}
	-\frac{2u(\zeta)}{(\zeta-w)^2}\dd\zeta = 0
	\iff \boxed{f'(w) = \frac 1{2\pi i}\int_{\abs\zeta = R}
	\frac{2 u(\zeta)}{(\zeta-w)^2}\dd\zeta}.\qedhere
    \end{equation*}
\end{soln}

\begin{problem}
    [p.174\#5]
    If $u(z)$ is harmonic and $0\le u(z)\le Ky$ for $y > 0$,
    prove that $u = ky$ with $0\le k\le K$.
\end{problem}

\begin{soln}
    Assume that $u$ is defined on the closed half plane $\ol\HH$.
    The given condition implies $u(z) = 0$ on the real line,
    so $u$ admits an extension $\tilde u\colon\CC\to\RR$
    with $\tilde u(\ol z) = -\tilde u(z)$ on all of $\CC$.
    Let $v$ be a conjugate harmonic function for $\tilde u$,
    so $f = \tilde u + iv$ is entire. By the previous problem,
    for a given $z$ and large enough $R$,
    \begin{equation*}
	f'(z) = \frac 1{2\pi i}\int_{\abs\zeta = R}
	\frac{2u(\zeta)}{(\zeta-z)^2}\dd\zeta.
    \end{equation*}
    By the given condition,
    $\abs{u(z)} \le K\cdot\operatorname{Im}z \le K\abs z$ for all $z$, so
    \begin{equation*}
	\abs{f'(z)} \le \frac 1{2\pi}\int_0^{2\pi}
	\abs{\frac{2u(Re^{i\theta})}{(Re^{i\theta} - z)^2}}\cdot R\dd\theta
	\le \frac 1{2\pi}\int_0^{2\pi}
	\frac{R^2K}{\abs{Re^{i\theta} - z}^2}\dd\theta
	= \frac 1{2\pi}\int_0^{2\pi}
	\frac{K}{\abs{e^{i\theta} - \frac zR}^2}\dd\theta.
    \end{equation*}
    For each $\eps$, there exists $R_\eps$ such that
    $\abs{e^{i\theta} - \frac zR}$ is within $\eps$ of $1$
    when $R > R_\eps$, so $\abs{f'(z)}$ is bounded above by $\frac K{1+\eps}$
    for sufficiently large $R$. By Liouville's theorem
    $f'$ is therefore constant, and more specifically we have bounded
    $\abs{f'(z)} \le K$.

    Finally, if $f(z) = az + b$ for constant $a$ and $b$,
    then the real part of $f(z)$ is zero whenever $z$ is real.
    This forces $a$ and $b$ to be purely imaginary:
    by taking $z = 0$, $b$ must have zero real part,
    and then taking $z = 1$ tells us that $a$ has zero real part.
    Thus $f(z) = (kz + \ell)i$ for $0\le k \le K$ and real $\ell$,
    so $\tilde u(z) = ky$ as desired.
\end{soln}











\end{document}
